\clearpage
\pdfbookmark[0]{Praefatio.}{praefatio}
\begin{center}\Large PRAEFATIO.\end{center}
\medskip
\SourceLeaf{010}Detectis legibus motus planetarum \textsc{Kepler}i ingenio non defuerunt subsidia ad singulorum planetarum elementa ex observationibus eruenda. \textsc{Tycho Brahe}, a quo astronomia practica ad fastigium antea ignotum evecta erat, cunctos planetas per longam annorum seriem summa cura tantaque perseverantia observaverat, ut \textsc{Kepler}o talis thesauri dignissimo heredi seligendi tantummodo cura restaret, quae ad scopum quemvis propositum facere viderentur. Nec mediocriter sublevabant hunc laborem motus planetarum medii summa iamdudum praecisione per observationes antiquissimas determinati.

Astronomi, qui post \textsc{Kepler}um conati sunt planetarum orbitas adiumento observationum recentiorum vel perfectiorum adhuc accuratius dimetiri, iisdem vel adhuc maioribus subsidiis adiuti sunt. Neque enim amplius de elementis plane incognitis eliciendis agebatur, sed nota leviter tantum corrigenda arctioribusque limitibus circumscribenda erant.

Principium gravitationis universalis a summo \textsc{Newton} detectum campum plane novum aperuit, legibusque iisdem, quibus quinque planetas regi \textsc{Kepler} expertus fuerat, levi tantum mutatione facta \emph{omnia} corpora coelestia necessario
\SourceLeaf{011}obsequi debere edocuit, quorum quidem motus a vi Solis tantum moderentur. Scilicet observationum testimonio fretus \textsc{Kepler} cuiusvis planetae orbitam ellipsem esse pronunciaverat, in qua areae circa Solem, focum alterum ellipsis occupantem, uniformiter describantur, et quidem ita, ut tempora revolutionum in ellipsibus diversis sint in ratione sesquialtera semiaxium maiorum. Contra \textsc{Newton}, principio gravitationis universalis posito, a priori demonstravit, corpora omnia a Solis vi attractiva gubernata in sectionibus conicis moveri debere, quarum quidem speciem unam, ellipses puta, planetae nobis exhibeant, dum species reliquae, parabolae et hyperbolae, pro aeque possibilibus haberi debeant, modo adsint corpora Solis vi velocitate debita occurrentia; Solem semper focum alterum sectionis conicae tenere; areas, quas corpus idem temporibus diversis circa Solem describat, his temporibus proportionales, areas denique a corporibus diversis, temporibus aequalibus, circa Solem descriptas, esse in ratione subduplicata semiparametrorum orbitarum: postrema harum legum, in motu elliptico cum ultima \textsc{Kepler}i lege identica, ad motum parabolicum hyperbolicumque patet, ad quos haecce applicari nequit, revolutionibus deficientibus. Iam filum repertum, quo ducente labyrinthum motuum cometarum antea inaccessum ingredi licuit. Quod tam feliciter successit, ut omnium cometarum motibus, qui quidem accurate observati essent, explicandis sufficeret unica hypothesis, orbitas parabolas esse. Ita systema gravitationis universalis novos analysi triumphos eosque splendidissimos paraverat; cometaeque usque ad illum diem semper indomiti, vel si devicti videbantur mox seditiosi et rebelles, frena sibi iniici passi, atque ex hostibus hospites redditi, iter suum in tramitibus a calculo delineatis prosequuti sunt, iisdem quibus planetae legibus aeternis religiose obtemperantes.

Iam in determinandis cometarum orbitis parabolicis ex observationibus difficultates suboriebantur longe maiores, quam in determinandis orbitis ellipticis planetarum, inde potissimum, quod cometae per brevius temporis 
\SourceLeaf{012}intervallum visi delectum observationum ad haec vel illa imprimis commodarum non concedebant, sed iis uti geometram cogebant, quas fors obtulerat, ita ut methodos speciales in calculis planetarum adhibitis vix unquam in usum vocare licuerit. Magnus ipse \textsc{Newton}, primus saeculi sui geometra, problematis difficultatem haud dissimulavit, attamen, ceu exspectari poterat, ex hoc quoque certamine victor evasit. Multi post \textsc{Newton}um geometrae eidem problemati operam suam navaverunt, varia utique fortuna, ita tamen, ut nostris temporibus parum desiderandum relictum sit.

Verum enim vero non est praetermittendum, in hoc quoque problemate peropportune difficultatem diminui per cognitionem unius elementi sectionis conicae, quum, per ipsam suppositionem orbitae parabolicae, axis maior infinite magnus statuatur. Quippe omnes parabolae, siquidem situs negligatur, per solam maiorem minoremve distantiam verticis a foco inter se differunt, dum sectiones conicae generaliter spectatae varietatem infinities maiorem admittant. Haud equidem aderat ratio sufficiens, cur cometarum traiectoriae absoluta praecisione parabolicae praesumerentur: quin potius infinite parum probabile censeri debet, rerum naturam unquam tali suppositioni annuisse. Attamen quum constaret, phaenomena corporis coelestis in ellipsi vel hyperbola incedentis, cuius axis maior permagnus sit ratione parametri, prope perihelium perparum discrepare a motu in parabola, cui eadem verticis a foco distantia, differentiamque eo leviorem evadere, quo maior fuerit illa ratio axis ad parametrum; porro quum experientia docuisset, inter motum observatum motumque in orbita parabolica computatum vix unquam maiores differentias remanere, quam quae ipsis observationum erroribus (hic plerumque satis notabilibus) tuto tribui poterant: astronomi apud parabolam subsistendum esse rati sunt. Recte sane, quum omnino deessent subsidia, e quibus, num ullae quantaeve differentiae a parabola adsint, satis certo colligi potuisset. Excipere oportet cometam celebrem \textsc{Halley}anum, qui ellipsem valde oblongam describens in
\SourceLeaf{013}reditu ad perihelium pluries observatus tempus periodicum nobis patefecit: tunc autem axi maiori inde cognito computus reliquorum elementorum vix pro difficiliori habendus est, quam determinatio orbitae parabolicae. Silentio quidem praeterire non possumus, astronomos etiam in nonnullis aliis cometis per tempus aliquanto longius observatis determinationem aberrationis a parabola tentavisse: attamen omnes methodi ad hunc finem propositae vel adhibitae innituntur suppositioni, discrepantiam a parabola haud considerabilem esse, quo pacto in illis tentaminibus ipsa parabola antea iam computata cognitionem approximatam singulorum elementorum (praeter axem maiorem vel tempus revolutionis inde pendens) iam subministravit, levibus tantum mutationibus corrigendam. Praeterea fatendum est, omnia ista tentamina vix unquam aliquid certi decidere valuisse, si forte cometam anni 1770 excipias.

Quamprimum motum planetae novi anno 1781 detecti cum hypothesi parabolica conciliari non posse cognitum est, astronomi orbitam circularem illi adaptare inchoaverunt, quod negotium per calculum perfacilem ac simplicem absolvitur. Fausta quadam fortuna orbita huius planetae mediocriter tantum excentrica erat, quo pacto elementa ·per suppositionem illam eruta saltem approximationem qualemcunque suppeditabant, cui dein determinationem elementorum ellipticorum superstruere licuit. Accedebant plura alia peropportuna. Quippe tardus planetae motus perparvaque orbitae ad planum eclipticae inclinatio non solum calculos longe simpliciores reddebant, methodosque speciales aliis casibus haud accommodandas in usum vocare concedebant, sed metum quoque dissipabant, ne planeta radiis Solis immersus postea quaeritantium curas eluderet (qui metus alias, praesertim si insuper lumen minus vividum fuisset, utique animos turbare potuisset), quo pacto accuratior orbitae determinatio tuto differri poterat, donec ex observationibus frequentioribus magisque remotis eligere liceret, quae ad propositum maxime commodae viderentur.

\SourceLeaf{014}In omnibus itaque casibus, ubi corporum coelestium orbitas ex observationibus deducere oportuit, commoda aderant quaedam haud spernenda methodorum specialium applicationem suadentia vel saltem permittentia, quorum commodorum potissimum id erat, quod per suppositiones hypotheticas cognitionem approximatam quorundam elementorum iamiam acquirere licuerat, antequam calculus elementorum ellipticorum susciperetur. Nihilominus satis mirum videtur, problema generale

\begin{quote}\itshape
Determinare orbitam corporis coelestis, absque omni suppositione hypothetica, ex observationibus tempus haud magnum complectentibus neque adeo delectum, pro applicatione methodorum specialium, patientibus
\end{quote}
\noindent usque ad initium huius saeculi penitus propemodum neglectum esse, vel saltem a nemine serio ac digne tractatum, quum certe theoreticis propter difficultatem atque elegantiam sese commendare potuisset, etiamsi apud practicos de summa eius utilitate nondum constaret. Scilicet invaluerat apud omnes opinio, \emph{impossibilem} esse talem determinationem completam ex observationibus breviori temporis intervallo inclusis, male sane fundata, quum nunc quidem certissimo iam evictum sit, orbitam corporis coelestis ex observationibus bonis paucos tantummodo dies complectentibus, absque ulla suppositione hypothetica, satis approximate iam determinari posse.

Incideram in quasdam ideas, quae ad solutionem problematis magni de quo dixi facere videbantur, mense Septembri a.\ 1801, tunc in labore plane diverso occupatus. Haud raro in tali casu, ne nimis a grata investigatione distrahamur, neglectas interire sinimus idearum associationes, quae attentius examinatae uberrimos fructus ferre potuissent. Forsan et illis ideolis eadem fortuna instabat, nisi peropportune incidissent in tempus, quo nullum sane faustius ad illas conservandas atque fovendas eligi potuisset. Scilicet eodem circiter tempore rumor de planeta novo Ian.\ 1 istius anni in specula Panormitana detecto per omnium ora volitabat, moxque ipsae observationes inde ab
\SourceLeaf{015}epocha illa usque ad 11 Febr. ab astronomo praestantissimo \textsc{Piazzi} institutae ad notitiam publicam pervenerunt. Nullibi sane in annalibus astronomiae occasionem tam gravem reperimus, vixque gravior excogitari posset, ad dignitatem istius problematis luculentissime ostendendam, quam tunc in tanto discrimine urgenteque necessitate, ubi omnis spes, atomum planetariam post annum fere elapsum in coelis inter innumeras stellulas reinveniendi, unice pendebat ab orbitae cognitione satis approximata, solis illis pauculis observationibus superstruenda. Unquamne opportunius experiri potuissem, ecquid valeant ideolae meae ad usum practicum, quam si tunc istis ad determinationem orbitae Cereris uterer, qui planeta inter 41 illos dies geocentrice arcum trium tantummodo graduum descripserat, et post annum elapsum in coeli plaga longissime illinc remota indagari debebat? Prima haecce methodi applicatio facta est mense Oct. 1801, primaque nox serena, ubi planeta ad normam numerorum inde deductorum quaesitus est\footnote{Dec. 7, 1801 a clar. de \textsc{Zach}.}, transfugam observationibus reddidit. Tres alii planetae novi inde ab illo tempore detecti occasiones novas suppeditaverunt methodi efficaciam ac generalitatem examinandi et comprobandi.

Optabant plures astronomi, statim post reinventionem Cereris, ut methodos ad istos calculos adhibitas publici iuris facerem; verum obstabant plura, quominus amicis hisce sollicitationibus tunc morem gererem: negotia alia, desiderium rem aliquando copiosius pertractandi, imprimisque expectatio, continuatam in hac disquisitione occupationem varias solutionis partes ad maius generalitatis, simplicitatis et elegantiae fastigium evecturam esse. Quae spes quum me haud fefellerit, non esse arbitror, cur me huius morae poeniteat. Methodi enim ab initio adhibitae identidem tot tantasque mutationes passae sunt, ut inter modum, quo olim orbita Cereris calculata est, institutionemque
\SourceLeaf{016}in hoc opere traditam vix ullum similitudinis vestigium remanserit. Quamquam vero a proposito meo alienum esset, de cunctis his disquisitionibus paullatim magis magisque perfectis narrationem completam perscribere, tamen in pluribus occasionibus, praesertim quoties de problemate quodam graviori agebatur, methodos anteriores quoque haud omnino supprimendas esse censui. Quin potius praeter problematum principalium solutiones plurima, quae in occupatione satis longa circa motus corporum coelestium in sectionibus conicis vel propter elegantiam analyticam vel imprimis propter usum practicum attentione digniora se mihi obtulerunt, in hoc opere exsequutus sum. Semper tamen vel rebus vel methodis mihi propriis maiorem curam dicavi, nota leviter tantum, quatenusque rerum nexus postulare videbatur, attingens.

Totum itaque opus in duas partes dividitur. In Libro primo evolvuntur relationes inter quantitates, a quibus motus corporum coelestium circa Solem secundum \textsc{Kepler}i leges pendet, et quidem in duabus primis Sectionibus relationes eae, ubi unicus tantum locus per se consideratur, in Sectione tertia et quarta vero eae, ubi plures loci inter se conferuntur. Illae continent expositionem methodorum tum vulgo usitatarum, tum potissimum aliarum illis ni fallor ad usum practicum longe praeferendarum, per quas ab elementis cognitis ad phaenomena descenditur; hae problemata multa gravissima tractant, quae viam ad operationes inversas sternunt. Scilicet quum ipsa phaenomena ex artificiosa intricataque quadam complicatione elementorum componantur, hanc texturae rationem penitius perspexisse oportet, antequam filorum explicationem operisque in elementa sua resolutionem cum spe successus suscipere liceat. Comparantur itaque in Libro primo instrumenta atque subsidia, per quae dein in Libro altero arduum hoc negotium ipsum perficitur: maxima laboris pars tunc iam in eo consistit, ut illa subsidia rite colligantur, ordine apto disponantur et in scopum propositum dirigantur.

Problemata graviora ad maximam partem per exempla idonea illustrata
\SourceLeaf{017}sunt, semper quoties quidem licuit ab observationibus non fictis desumta. Ita non solum methodorum efficaciae maior fiducia conciliabitur, ususque clarius ob oculos ponetur, sed id quoque cautum iri spero, ut nec minus exercitati a studio harum rerum deterreantur, quae procul dubio partem foecundissimam et pulcherrimam astronomiae theoricae constituunt.

Scripsi Gottingae d. 28 Martii 1809.
